\chapter{系统硬件设计}

\section{硬件设计概述}

本节填写系统硬件设计的总体说明。可简要介绍本课程设计所使用的硬件平台、主控芯片或核心器件、输入输出模块、显示模块、通信模块、电源模块以及下载调试接口等内容。

硬件设计章节应重点说明系统由哪些硬件部分组成，各硬件模块之间如何连接，以及这些硬件资源如何支撑系统功能实现。本节不需要展开过多程序实现细节，软件流程和核心代码可放在后续软件设计章节中说明。

撰写时可按照“主控平台—输入模块—输出模块—显示模块—通信调试模块—电源与下载接口”的逻辑展开，也可以根据实际课题特点调整章节结构。

\section{主控最小系统}

本节填写主控芯片或核心控制器的最小系统设计内容。主控最小系统通常包括主控芯片、电源电路、复位电路、时钟电路、下载调试接口以及必要的启动配置电路等。

其中，电源电路用于为主控芯片及外围模块提供稳定工作电压；复位电路用于在上电或异常情况下使系统恢复到初始状态；时钟电路用于为主控芯片和片上外设提供工作时序；下载调试接口用于程序下载、在线调试和实验验证。

在实际撰写时，应结合所用开发板或自制电路，说明主控芯片型号、主要资源特点、供电方式、复位方式、时钟来源和下载调试方式。若课程设计基于现成开发板完成，也可说明本设计主要使用开发板上已经集成的最小系统资源。

\section{输入模块设计}

本节填写系统输入模块的硬件设计内容。输入模块可以是按键、传感器、开关量输入、模拟量采集模块、通信输入模块或其他人机交互输入资源。

输入模块设计应说明输入信号来源、连接接口、信号类型和基本工作原理。例如，数字输入模块可说明高低电平变化与系统状态之间的关系；模拟输入模块可说明采样通道、信号范围和转换方式；传感器模块可说明检测对象、输出信号和主控读取方式。

如果输入模块涉及消抖、滤波、阈值判断、上拉下拉配置、外部中断或定时采样等内容，可在本节进行硬件层面的简要说明，具体的软件判断逻辑可放在第四章展开。

\section{输出与显示模块设计}

本节填写系统输出模块和显示模块的硬件设计内容。输出模块可以包括 LED 指示灯、蜂鸣器、继电器、电机、数码管、LCD、OLED 或其他执行与显示设备。

输出模块设计应说明输出设备的连接方式、控制信号、工作电平和功能作用。例如，LED 指示灯可用于显示系统运行状态或事件反馈；显示模块可用于显示系统参数、状态信息、测试结果或调试信息；执行模块可用于完成具体控制动作。

若系统使用显示屏，应说明显示模块的接口方式、显示内容和在系统测试中的作用。若系统未使用显示模块，可将本节改为“执行模块设计”或“状态指示模块设计”。

\section{通信与调试接口设计}

本节填写系统通信接口和调试接口设计内容。常见通信接口包括串口、I\textsuperscript{2}C、SPI、CAN、USB、以太网或无线通信模块等；常见调试接口包括 JTAG、SWD、ISP、下载器接口或仿真器接口等。

通信接口部分可说明通信模块的用途、连接方式、通信参数和数据交互内容。例如，可用于调试输出、上位机通信、模块配置、数据上传或远程控制等。

调试接口部分可说明程序下载、在线调试和实验验证所使用的接口。若课程设计基于开发板完成，可说明开发板已提供相应调试接口，本设计主要利用其完成程序下载和运行测试。

\section{硬件连接与资源分配}

本节填写系统主要硬件资源、连接接口和功能说明。为了使硬件设计更加清晰，建议在本节加入硬件资源与接口分配表，如表~\ref{tab:hardware_resources} 所示。

\begin{table}[htbp]
    \centering
    \caption{系统主要硬件资源与接口分配}
    \label{tab:hardware_resources}
    \small
    \setlength{\tabcolsep}{4pt}
    \renewcommand{\arraystretch}{1.25}
    \begin{tabular}{|p{0.18\textwidth}|p{0.28\textwidth}|p{0.44\textwidth}|}
        \hline
        硬件资源 & 连接/接口 & 功能说明 \\
        \hline
        主控芯片/核心板 & 填写芯片型号或核心板接口 & 作为系统控制核心，完成输入采集、逻辑处理、输出控制和状态管理。 \\
        \hline
        输入模块 & 填写引脚、接口或通信方式 & 用于采集外部输入信号，为系统控制逻辑提供判断依据。 \\
        \hline
        输出模块 & 填写引脚、接口或驱动方式 & 用于输出控制结果、状态提示或执行动作。 \\
        \hline
        显示模块 & 填写显示接口或连接方式 & 用于显示系统状态、参数信息、测试结果或调试信息。 \\
        \hline
        通信模块 & 填写通信接口和主要参数 & 用于调试输出、数据通信或上位机交互。 \\
        \hline
        下载调试接口 & 填写下载或调试接口 & 用于程序下载、在线调试和实验验证。 \\
        \hline
        电源模块 & 填写供电方式和电压范围 & 为主控芯片及外围模块提供稳定工作电源。 \\
        \hline
    \end{tabular}
\end{table}

实际写作时，应将表~\ref{tab:hardware_resources} 中的“填写……”内容替换为本课程设计所使用的真实硬件资源和接口信息。若某些模块未使用，可删除对应行；若系统包含更多外设，也可继续增加表格行。

\section{本章小结}

本章主要介绍了课程设计系统的硬件设计内容，包括硬件总体结构、主控最小系统、输入模块、输出与显示模块、通信调试接口以及硬件资源分配等。通过本章设计，可以明确系统硬件组成和接口关系，为后续软件设计、系统调试和功能测试奠定基础。